\subsection*{МГУ-13.8}
\setcounter{equation}{0}

\begin{abstract}
Длинный прямолинейный провод радиусом $а$ из проводящего материала с магнитной проницаемостью $\mu$ находится в воздухе. По проводу течет постоянный ток $I$, равномерно распределенный по его сечению. Найти намагниченность $М$ провода на расстоянии $r$ от его оси, объемную $j_m$ и поверхностную $i_m$ плотность токов намагничивания.
\end{abstract}

\noindent \hrulefill
\\

Согласно теореме Ампера для окружностей радиуса $r$:

$$H_{(r)}2 \pi r = I \frac{\pi r^2}{\pi a^2} \xrightarrow[]{} H_{(r)} = \frac{Ir}{2 \pi a^2}$$

Намагниченность провода на расстоянии $r$ от его оси: 

$$\vec{B} = \mu_0 (\vec{H} + \vec{M}) = \mu \mu_0 \vec{H}$$

$$ \xrightarrow[]{} M_{(r)} = (\mu -1) H_{(r)} = (\mu -1) \frac{Ir}{2 \pi a^2} $$

Объемная $j_m$ плотность токов намагничивания:

$$\vec{j_m} = \frac{1}{r}.\frac{d(r.M)}{dr} \hat{z} = \frac{1}{r}.\frac{d}{dr} [(\mu -1) \frac{Ir^2}{2 \pi a^2}] \hat{z} = \frac{(\mu - 1) I}{\pi a^2} \hat{z}$$

Поверхностную $i_m$ плотность токов намагничивания:

$$\vec{i_m} = \vec{M} \times \vec{n} = M_{(a)}(\vec{\phi} \times \hat{r}) = -M_{(a)} \hat{z}$$

$$M_{(a)} = (\mu -1) \frac{I}{2 \pi a} \xrightarrow[]{} \vec{i_m} = - \frac{(\mu -1)I}{2 \pi a} \hat{z}$$

\textbf{Ответ:}

$$ M_{(r)} = (\mu -1) \frac{Ir}{2 \pi a^2}$$

$$\vec{j_m} = \frac{(\mu - 1) I}{\pi a^2} \hat{z}$$

$$\vec{i_m} = - \frac{(\mu -1)I}{2 \pi a} \hat{z}$$







